% preambulo-simbolos.tex
% Fallback de fonte para símbolos que a fonte principal não tem.
% Instalar em: <userdir>\preambulo-simbolos.tex
% Usado por academico.yaml, caderno.yaml e caderno-frente.yaml.
%
% POR QUE ISTO EXISTE
% XeLaTeX não faz fallback automático de fonte. Se a fonte principal não
% tem o glifo, o caractere some do PDF e o único sinal é um
% "[WARNING] Missing character" na saída. Palatino Linotype cobre latim e
% grego politônico, mas não tem setas, subscritos nem colchetes angulares.
%
% Cada \newunicodechar abaixo desvia UM caractere para a fonte de símbolos.
%
% COMO ACRESCENTAR UM CARACTERE
% Copie uma linha, troque o caractere e o código hexadecimal:
%     \newunicodechar{X}{\FBsym{"XXXX}}
% O código aparece no próprio aviso do Pandoc: "no → (U+2192)" -> "2192.
%
% CUIDADO COM A SINTAXE
% Escrever \newcommand\FB[1]{{\symbolfont #1}} e passar o caractere como
% argumento QUEBRA o build, com "Argument of \MT@is@char has an extra }"
% — é conflito com o microtype. Por isso aqui se usa \symbol{"código}
% dentro de \bgroup...\egroup, que não passa o caractere adiante.

\usepackage{newunicodechar}

% Fonte de fallback para os símbolos que a Cardo não tem.
% A Cardo é excelente em latim, grego politônico e sinais filológicos,
% mas não traz setas, subscritos nem colchetes angulares matemáticos.
% "Segoe UI Symbol" vem com o Windows e cobre tudo o que está mapeado
% abaixo. Alternativas testadas: DejaVu Sans, Segoe UI, Cambria Math.
% Se o build reclamar que a fonte não existe, troque só esta linha.
\newfontfamily\symbolfont{Segoe UI Symbol}[Scale=MatchLowercase]

\newcommand\FBsym[1]{\bgroup\symbolfont\symbol{#1}\egroup}

% --- setas (aparecem nos handouts, em diagramas) ---
\newunicodechar{→}{\FBsym{"2192}}
\newunicodechar{←}{\FBsym{"2190}}
\newunicodechar{↑}{\FBsym{"2191}}
\newunicodechar{↓}{\FBsym{"2193}}
\newunicodechar{↔}{\FBsym{"2194}}
\newunicodechar{⇒}{\FBsym{"21D2}}
\newunicodechar{⇔}{\FBsym{"21D4}}

% --- colchetes angulares (inserção editorial: ⟨texto do editor⟩) ---
\newunicodechar{⟨}{\FBsym{"27E8}}
\newunicodechar{⟩}{\FBsym{"27E9}}

% --- subscritos (Barnes usa em fórmulas) ---
\newunicodechar{₀}{\FBsym{"2080}}
\newunicodechar{₁}{\FBsym{"2081}}
\newunicodechar{₂}{\FBsym{"2082}}
\newunicodechar{₃}{\FBsym{"2083}}
\newunicodechar{₄}{\FBsym{"2084}}

% --- matemáticos e sinais diversos ---
\newunicodechar{≥}{\FBsym{"2265}}
\newunicodechar{≤}{\FBsym{"2264}}
\newunicodechar{≠}{\FBsym{"2260}}
\newunicodechar{≈}{\FBsym{"2248}}
\newunicodechar{∴}{\FBsym{"2234}}
% ⸫ (U+2E2B) é o "portanto" que aparece nas reconstruções de argumento dos
% handouts. Quase nenhuma fonte o tem — aqui ele é desenhado com o ∴ (U+2234),
% que é o símbolo lógico correto e existe em toda parte.
\newunicodechar{⸫}{\FBsym{"2234}}
\newunicodechar{∙}{\FBsym{"2219}}
\newunicodechar{★}{\FBsym{"2605}}

% --- latim estendido raro (transliteração do grego) ---
\newunicodechar{ḗ}{\FBsym{"1E17}}
\newunicodechar{ṓ}{\FBsym{"1E53}}
